\section{Preliminary}
The extended real-valued function $f: \mathbb{R}^{q}\rightarrow (-\infty, +\infty]$ is called \textit{proper} if its domain $\text{dom}(f) := \{\bx \in \calX\mid f(\bx) < +\infty\}$ is nonempty and $\inf\limits_{\bx \in \mathbb{R}^q}f > -\infty$; it is \textit{closed} if its epigraph $\text{epi}(f) := \left\{ (\bx, t) \in \mathbb{R}^{q+1}\mid f(\bx) \leq t \right\}$ is closed.  For a proper, closed, and convex function $f$, the \textit{proximal mapping} of $f$, denoted by $\operatorname{prox}_{\tau f}(\by)$, is defined as
\begin{align}
    % \label{Def: the definition of proximal operator}
    \operatorname{prox}_{\tau f}(\by) := \argmin\limits_{\bx \in \operatorname{dom}(f)}\left\{ f(\bx) + \frac{1}{2\tau}\|\bx - \by\|^{2}\right\}, \notag
\end{align} 
where $\tau$ is a positive scalar parameter. The proximal mapping admits the following
properties.
\begin{proposition}[\cite{Rockafellar1998VariationalA}]
    Let $f : \mathbb{R}^{q}\rightarrow \mathbb{R}\cup \{+\infty\}$ be proper and
    closed. Then, the following holds.
    \begin{enumerate}[label=(\roman{*})]
        \item The set $\operatorname{prox}_{\tau f}(\bx)$ is nonempty and compact
            for any $\bx \in \mathbb{R}^{q}$ and $\tau > 0$.

        \item If $f$ is convex, then $\operatorname{prox}_{\tau f}(\bx)$
            contains exactly one value for any $\bx \in \mathbb{R}^{q}$ and
            $\tau > 0$.
    \end{enumerate}
\end{proposition}

Let us introduce the StoRM (Stochastic Recursive Momentum) variance reduction estimator \cite{cutkoskyMomentumBasedVarianceReduction2019}. In stochastic optimization problems where $F(\bx) = \mathbb{E}_{\xi}[F(\bx;\xi)]$ with $\xi \in \Xi$, 
% standard momentum methods use
% \begin{align*}
%     \bd^{k}= \rho_{k}\nabla F(\bx^{k}, \xi^{k}) + (1 - \rho_{k}) \bd^{k-1}, \text{ where }\rho_{k}\in [0, 1],
% \end{align*}
% yielding a biased estimate of $\nabla F(\bx^{k})$.
StoRM adds a
correction term, $(1 - \rho_{k}) (\nabla F(\bx^{k}, \xi^{k}) - \nabla F(\bx^{k-1}, \xi^{k}))$, resulting in:
\begin{align*}
     & \bd^{k}=\left\{\begin{array}{ll}\nabla F\left(\bx^{k}, \xi^{k}\right), & \text{if } k=1, \\ \nabla F\left(\bx^{k}, \xi^{k}\right) + \left(1-\rho_k\right)\left( \bd^{k-1} - \nabla F\left(\bx^{k-1}, \xi^{k}\right) \right), & \text{if } k>1,\end{array}\right. \\
     & \bx^{k+1}=\bx^{k}-\eta_{k}\bd^{k}.
\end{align*}
This exploits the smoothness of $F(\cdot, \xi)$, leading to efficient
variance reduction. When $\rho_{k}= 1$, $\bd^{k}$ reduces to vanilla mini-batch SGD. We set $\rho_{k}\in (0,1)$ for $k \geq 1$ in this paper

A function $f: \calX \rightarrow \RR$ is said to be $L_f$-smooth if  $\| \nabla f(\bx) - \nabla f(\bx') \| \leq L_f \| \bx - \bx'\|$ for all $\bx,\, \bx' \in \RR^n$. A mapping $\mathcal{M} : \mathcal{D} \to \RR^l $ is said to be $C$-Lipschitz continuous if $\|\mathcal{M}(\bx) - \mathcal{M}(\bx')\| \leq C\|\bx - \bx'\|$ for all $\bx, \bx' \in \mathcal{D}$. 

% A well-known gradient descent lemma for a smooth function is that
% \begin{align}
% % \label{Eq: the descent lemma}
%     f(\bx) \leq f(\bx') + \langle \nabla f(\bx'), \bx - \bx' \rangle + \frac{L_f}{2} \| \bx - \bx'\|^2. \notag
% \end{align}
